\clearpage\phantomsection\label{couverture-programme}\pdfbookmark[0]{Repérage des attendus}{bm-couverture-programme}\gdef\otChapNum{0}\pagestyle{otFront}\otPageTitre{Repérage des attendus}{Quatrième 2026-2027 · Lecture par notions}
Cette matrice relie les connaissances des attendus annuels et les acquis
réactivés à des exercices précis. Les quatre repères de chaque chapitre
permettent de retrouver un travail ciblé. Les séries indiquées ne sont pas
exhaustives ; les cours et les défis complètent ce parcours.\par\bigskip
Points à retrouver : préfixes de nano à giga (3), carrés jusqu’à 144 et
racines en contexte (4), équations du premier degré (6), grandeurs composées
(8), graphiques sans formalisme de fonctions (9), médiane sur de petits
effectifs et diagrammes circulaires (10), événement contraire (11),
cas d’égalité des triangles (12), Pythagore et réciproque (13),
Thalès emboîté et réciproque (14), cosinus (15), translations et effets
sur les grandeurs (16), patrons et repérage dans le pavé (17),
programmation de niveaux 1 et 2 (18).\par\bigskip
\begin{attention}Les difficultés signalées par trois pastilles et les défis sont des choix d’entraînement. Ils ne changent pas le programme officiel et ne constituent pas une liste d’exigences supplémentaires pour un examen.\end{attention}
\clearpage
\otPageTitre{Repères et exercices}{Chapitres 1 à 3}\fontsize{10}{13}\selectfont
\otIntertitre{1. Calculer avec les nombres relatifs}\begin{tabularx}{\linewidth}{@{}Yl@{}}\toprule
Produits et quotients de relatifs & \hyperref[exo-1-1]{1.1}, \hyperref[exo-1-2]{1.2}, \hyperref[exo-1-3]{1.3}, \hyperref[exo-1-5]{1.5}\\[2mm]
Opposés et règles des signes & \hyperref[exo-1-4]{1.4}, \hyperref[exo-1-6]{1.6}, \hyperref[exo-1-14]{1.14}\\[2mm]
Priorités et calcul réfléchi & \hyperref[exo-1-7]{1.7}, \hyperref[exo-1-8]{1.8}, \hyperref[exo-1-9]{1.9}, \hyperref[exo-1-10]{1.10}, \hyperref[exo-1-15]{1.15}\\[2mm]
Modélisation et contrôle & \hyperref[exo-1-11]{1.11}, \hyperref[exo-1-12]{1.12}, \hyperref[exo-1-13]{1.13}\\[2mm]
\bottomrule\end{tabularx}\par\vspace{4mm}
\otIntertitre{2. Fractions et nombres premiers}\begin{tabularx}{\linewidth}{@{}Yl@{}}\toprule
Fractions égales et comparaison & \hyperref[exo-2-1]{2.1}, \hyperref[exo-2-2]{2.2}, \hyperref[exo-2-10]{2.10}\\[2mm]
Quatre opérations et inverse & \hyperref[exo-2-3]{2.3}, \hyperref[exo-2-4]{2.4}, \hyperref[exo-2-5]{2.5}, \hyperref[exo-2-8]{2.8}, \hyperref[exo-2-9]{2.9}\\[2mm]
Nombres premiers et simplification & \hyperref[exo-2-6]{2.6}, \hyperref[exo-2-7]{2.7}, \hyperref[exo-2-12]{2.12}\\[2mm]
Partages et fractions d’un reste & \hyperref[exo-2-11]{2.11}, \hyperref[exo-2-13]{2.13}, \hyperref[exo-2-14]{2.14}, \hyperref[exo-2-15]{2.15}\\[2mm]
\bottomrule\end{tabularx}\par\vspace{4mm}
\otIntertitre{3. Puissances et notation scientifique}\begin{tabularx}{\linewidth}{@{}Yl@{}}\toprule
Puissances d’exposant positif & \hyperref[exo-3-1]{3.1}, \hyperref[exo-3-2]{3.2}, \hyperref[exo-3-14]{3.14}\\[2mm]
Puissances de dix positives et négatives & \hyperref[exo-3-3]{3.3}, \hyperref[exo-3-7]{3.7}, \hyperref[exo-3-9]{3.9}, \hyperref[exo-3-10]{3.10}\\[2mm]
Notation scientifique et comparaison & \hyperref[exo-3-4]{3.4}, \hyperref[exo-3-6]{3.6}, \hyperref[exo-3-8]{3.8}, \hyperref[exo-3-13]{3.13}, \hyperref[exo-3-15]{3.15}\\[2mm]
Préfixes de nano à giga & \hyperref[exo-3-5]{3.5}, \hyperref[exo-3-11]{3.11}, \hyperref[exo-3-12]{3.12}\\[2mm]
\bottomrule\end{tabularx}\par\vspace{4mm}
\clearpage
\otPageTitre{Repères et exercices}{Chapitres 4 à 6}\fontsize{10}{13}\selectfont
\otIntertitre{4. Carrés et racines carrées}\begin{tabularx}{\linewidth}{@{}Yl@{}}\toprule
Carrés parfaits de 1 à 144 & \hyperref[exo-4-1]{4.1}, \hyperref[exo-4-2]{4.2}, \hyperref[exo-4-3]{4.3}\\[2mm]
Définition de la racine carrée & \hyperref[exo-4-4]{4.4}, \hyperref[exo-4-5]{4.5}, \hyperref[exo-4-8]{4.8}\\[2mm]
Encadrements et valeurs approchées & \hyperref[exo-4-6]{4.6}, \hyperref[exo-4-10]{4.10}, \hyperref[exo-4-14]{4.14}\\[2mm]
Aires et longueurs exactes & \hyperref[exo-4-7]{4.7}, \hyperref[exo-4-9]{4.9}, \hyperref[exo-4-11]{4.11}, \hyperref[exo-4-12]{4.12}, \hyperref[exo-4-13]{4.13}, \hyperref[exo-4-15]{4.15}\\[2mm]
\bottomrule\end{tabularx}\par\vspace{4mm}
\otIntertitre{5. Expressions et distributivité}\begin{tabularx}{\linewidth}{@{}Yl@{}}\toprule
Structure, substitution et réduction & \hyperref[exo-5-1]{5.1}, \hyperref[exo-5-2]{5.2}, \hyperref[exo-5-3]{5.3}, \hyperref[exo-5-10]{5.10}\\[2mm]
Développer par distributivité simple & \hyperref[exo-5-4]{5.4}, \hyperref[exo-5-6]{5.6}, \hyperref[exo-5-7]{5.7}\\[2mm]
Factoriser avec un facteur commun & \hyperref[exo-5-5]{5.5}, \hyperref[exo-5-8]{5.8}, \hyperref[exo-5-13]{5.13}\\[2mm]
Programmes de calcul et preuves & \hyperref[exo-5-9]{5.9}, \hyperref[exo-5-11]{5.11}, \hyperref[exo-5-12]{5.12}, \hyperref[exo-5-14]{5.14}, \hyperref[exo-5-15]{5.15}\\[2mm]
\bottomrule\end{tabularx}\par\vspace{4mm}
\otIntertitre{6. Équations et résolution de problèmes}\begin{tabularx}{\linewidth}{@{}Yl@{}}\toprule
Inconnue et test d’une solution & \hyperref[exo-6-1]{6.1}, \hyperref[exo-6-5]{6.5}\\[2mm]
Équations du premier degré & \hyperref[exo-6-2]{6.2}, \hyperref[exo-6-3]{6.3}, \hyperref[exo-6-4]{6.4}, \hyperref[exo-6-6]{6.6}, \hyperref[exo-6-7]{6.7}, \hyperref[exo-6-8]{6.8}\\[2mm]
Mise en équation et vérification & \hyperref[exo-6-11]{6.11}, \hyperref[exo-6-12]{6.12}, \hyperref[exo-6-13]{6.13}, \hyperref[exo-6-14]{6.14}, \hyperref[exo-6-15]{6.15}\\[2mm]
Interprétation et cas particuliers & \hyperref[exo-6-9]{6.9}, \hyperref[exo-6-10]{6.10}\\[2mm]
\bottomrule\end{tabularx}\par\vspace{4mm}
\clearpage
\otPageTitre{Repères et exercices}{Chapitres 7 à 9}\fontsize{10}{13}\selectfont
\otIntertitre{7. Proportionnalité et pourcentages}\begin{tabularx}{\linewidth}{@{}Yl@{}}\toprule
Quatrième proportionnelle et méthodes & \hyperref[exo-7-1]{7.1}, \hyperref[exo-7-2]{7.2}, \hyperref[exo-7-6]{7.6}\\[2mm]
Pourcentages directs et inverses & \hyperref[exo-7-3]{7.3}, \hyperref[exo-7-4]{7.4}, \hyperref[exo-7-7]{7.7}, \hyperref[exo-7-11]{7.11}, \hyperref[exo-7-13]{7.13}, \hyperref[exo-7-15]{7.15}\\[2mm]
Échelles et formules proportionnelles & \hyperref[exo-7-5]{7.5}, \hyperref[exo-7-10]{7.10}, \hyperref[exo-7-12]{7.12}, \hyperref[exo-7-14]{7.14}\\[2mm]
Graphiques et proportionnalité & \hyperref[exo-7-8]{7.8}, \hyperref[exo-7-9]{7.9}\\[2mm]
\bottomrule\end{tabularx}\par\vspace{4mm}
\otIntertitre{8. Grandeurs composées et unités}\begin{tabularx}{\linewidth}{@{}Yl@{}}\toprule
Vitesses, durées et distances & \hyperref[exo-8-1]{8.1}, \hyperref[exo-8-2]{8.2}, \hyperref[exo-8-3]{8.3}, \hyperref[exo-8-6]{8.6}, \hyperref[exo-8-11]{8.11}, \hyperref[exo-8-14]{8.14}\\[2mm]
Débits et conversions composées & \hyperref[exo-8-4]{8.4}, \hyperref[exo-8-5]{8.5}, \hyperref[exo-8-7]{8.7}, \hyperref[exo-8-8]{8.8}, \hyperref[exo-8-13]{8.13}, \hyperref[exo-8-15]{8.15}\\[2mm]
Concentration et masse volumique & \hyperref[exo-8-9]{8.9}, \hyperref[exo-8-12]{8.12}\\[2mm]
Grandeurs produits et cohérence & \hyperref[exo-8-10]{8.10}\\[2mm]
\bottomrule\end{tabularx}\par\vspace{4mm}
\otIntertitre{9. Relier des grandeurs et lire des graphiques}\begin{tabularx}{\linewidth}{@{}Yl@{}}\toprule
Formules et tableaux de valeurs & \hyperref[exo-9-1]{9.1}, \hyperref[exo-9-2]{9.2}, \hyperref[exo-9-9]{9.9}\\[2mm]
Lecture graphique dans les deux sens & \hyperref[exo-9-3]{9.3}, \hyperref[exo-9-4]{9.4}, \hyperref[exo-9-7]{9.7}\\[2mm]
Construction et proportionnalité & \hyperref[exo-9-5]{9.5}, \hyperref[exo-9-6]{9.6}, \hyperref[exo-9-8]{9.8}\\[2mm]
Interprétation, tableur et domaine & \hyperref[exo-9-10]{9.10}, \hyperref[exo-9-11]{9.11}, \hyperref[exo-9-12]{9.12}, \hyperref[exo-9-13]{9.13}, \hyperref[exo-9-14]{9.14}, \hyperref[exo-9-15]{9.15}\\[2mm]
\bottomrule\end{tabularx}\par\vspace{4mm}
\clearpage
\otPageTitre{Repères et exercices}{Chapitres 10 à 12}\fontsize{10}{13}\selectfont
\otIntertitre{10. Statistiques et diagrammes}\begin{tabularx}{\linewidth}{@{}Yl@{}}\toprule
Effectifs, fréquences et moyenne & \hyperref[exo-10-1]{10.1}, \hyperref[exo-10-2]{10.2}, \hyperref[exo-10-6]{10.6}, \hyperref[exo-10-11]{10.11}\\[2mm]
Médiane sur un effectif au plus égal à 30 & \hyperref[exo-10-3]{10.3}, \hyperref[exo-10-4]{10.4}, \hyperref[exo-10-7]{10.7}, \hyperref[exo-10-12]{10.12}, \hyperref[exo-10-15]{10.15}\\[2mm]
Diagrammes circulaires et en bâtons & \hyperref[exo-10-5]{10.5}, \hyperref[exo-10-8]{10.8}, \hyperref[exo-10-9]{10.9}\\[2mm]
Tableur et interprétation critique & \hyperref[exo-10-10]{10.10}, \hyperref[exo-10-13]{10.13}, \hyperref[exo-10-14]{10.14}\\[2mm]
\bottomrule\end{tabularx}\par\vspace{4mm}
\otIntertitre{11. Probabilités : comprendre le hasard}\begin{tabularx}{\linewidth}{@{}Yl@{}}\toprule
Vocabulaire et probabilités entre 0 et 1 & \hyperref[exo-11-1]{11.1}, \hyperref[exo-11-2]{11.2}, \hyperref[exo-11-3]{11.3}\\[2mm]
Équiprobabilité et écritures diverses & \hyperref[exo-11-4]{11.4}, \hyperref[exo-11-6]{11.6}, \hyperref[exo-11-7]{11.7}\\[2mm]
Événements contraires & \hyperref[exo-11-5]{11.5}, \hyperref[exo-11-8]{11.8}, \hyperref[exo-11-11]{11.11}, \hyperref[exo-11-13]{11.13}\\[2mm]
Modèles non équiprobables et aires & \hyperref[exo-11-9]{11.9}, \hyperref[exo-11-10]{11.10}, \hyperref[exo-11-12]{11.12}, \hyperref[exo-11-14]{11.14}, \hyperref[exo-11-15]{11.15}\\[2mm]
\bottomrule\end{tabularx}\par\vspace{4mm}
\otIntertitre{12. Construire et comparer des triangles}\begin{tabularx}{\linewidth}{@{}Yl@{}}\toprule
Existence et protocoles de construction & \hyperref[exo-12-1]{12.1}, \hyperref[exo-12-2]{12.2}, \hyperref[exo-12-3]{12.3}, \hyperref[exo-12-4]{12.4}, \hyperref[exo-12-14]{12.14}\\[2mm]
Trois cas d’égalité des triangles & \hyperref[exo-12-5]{12.5}, \hyperref[exo-12-6]{12.6}, \hyperref[exo-12-7]{12.7}, \hyperref[exo-12-8]{12.8}\\[2mm]
Correspondance de côtés et d’angles & \hyperref[exo-12-9]{12.9}, \hyperref[exo-12-10]{12.10}, \hyperref[exo-12-11]{12.11}, \hyperref[exo-12-13]{12.13}\\[2mm]
Raisonnements et preuves géométriques & \hyperref[exo-12-12]{12.12}, \hyperref[exo-12-15]{12.15}\\[2mm]
\bottomrule\end{tabularx}\par\vspace{4mm}
\clearpage
\otPageTitre{Repères et exercices}{Chapitres 13 à 15}\fontsize{10}{13}\selectfont
\otIntertitre{13. Le théorème de Pythagore}\begin{tabularx}{\linewidth}{@{}Yl@{}}\toprule
Calculer une hypoténuse & \hyperref[exo-13-1]{13.1}, \hyperref[exo-13-3]{13.3}, \hyperref[exo-13-7]{13.7}, \hyperref[exo-13-9]{13.9}\\[2mm]
Calculer un côté de l’angle droit & \hyperref[exo-13-2]{13.2}, \hyperref[exo-13-4]{13.4}, \hyperref[exo-13-8]{13.8}, \hyperref[exo-13-12]{13.12}\\[2mm]
Réciproque et triangle non rectangle & \hyperref[exo-13-5]{13.5}, \hyperref[exo-13-6]{13.6}, \hyperref[exo-13-10]{13.10}\\[2mm]
Applications et précision des résultats & \hyperref[exo-13-11]{13.11}, \hyperref[exo-13-13]{13.13}, \hyperref[exo-13-14]{13.14}, \hyperref[exo-13-15]{13.15}\\[2mm]
\bottomrule\end{tabularx}\par\vspace{4mm}
\otIntertitre{14. Thalès : triangles emboîtés}\begin{tabularx}{\linewidth}{@{}Yl@{}}\toprule
Thalès dans des triangles emboîtés & \hyperref[exo-14-1]{14.1}, \hyperref[exo-14-2]{14.2}, \hyperref[exo-14-3]{14.3}, \hyperref[exo-14-7]{14.7}, \hyperref[exo-14-8]{14.8}, \hyperref[exo-14-10]{14.10}, \hyperref[exo-14-12]{14.12}\\[2mm]
Réciproque et absence de parallélisme & \hyperref[exo-14-4]{14.4}, \hyperref[exo-14-5]{14.5}, \hyperref[exo-14-6]{14.6}, \hyperref[exo-14-13]{14.13}\\[2mm]
Constructions et partage de segments & \hyperref[exo-14-9]{14.9}, \hyperref[exo-14-14]{14.14}\\[2mm]
Problèmes de longueurs et d’aires & \hyperref[exo-14-11]{14.11}, \hyperref[exo-14-15]{14.15}\\[2mm]
\bottomrule\end{tabularx}\par\vspace{4mm}
\otIntertitre{15. Le cosinus dans le triangle rectangle}\begin{tabularx}{\linewidth}{@{}Yl@{}}\toprule
Hypoténuse et côté adjacent & \hyperref[exo-15-1]{15.1}, \hyperref[exo-15-2]{15.2}, \hyperref[exo-15-3]{15.3}, \hyperref[exo-15-10]{15.10}\\[2mm]
Calculer une longueur par le cosinus & \hyperref[exo-15-4]{15.4}, \hyperref[exo-15-5]{15.5}, \hyperref[exo-15-6]{15.6}, \hyperref[exo-15-9]{15.9}\\[2mm]
Déterminer un angle en degrés & \hyperref[exo-15-7]{15.7}, \hyperref[exo-15-8]{15.8}\\[2mm]
Combiner cosinus et Pythagore & \hyperref[exo-15-11]{15.11}, \hyperref[exo-15-12]{15.12}, \hyperref[exo-15-13]{15.13}, \hyperref[exo-15-14]{15.14}, \hyperref[exo-15-15]{15.15}\\[2mm]
\bottomrule\end{tabularx}\par\vspace{4mm}
\clearpage
\otPageTitre{Repères et exercices}{Chapitres 16 à 18}\fontsize{10}{13}\selectfont
\otIntertitre{16. Translations, agrandissements et réductions}\begin{tabularx}{\linewidth}{@{}Yl@{}}\toprule
Construire par translation & \hyperref[exo-16-1]{16.1}, \hyperref[exo-16-2]{16.2}, \hyperref[exo-16-3]{16.3}, \hyperref[exo-16-13]{16.13}\\[2mm]
Conservation et parallélogrammes & \hyperref[exo-16-4]{16.4}, \hyperref[exo-16-6]{16.6}\\[2mm]
Frises et pavages & \hyperref[exo-16-5]{16.5}, \hyperref[exo-16-15]{16.15}\\[2mm]
Rapports de longueurs, d’aires et de volumes & \hyperref[exo-16-7]{16.7}, \hyperref[exo-16-8]{16.8}, \hyperref[exo-16-9]{16.9}, \hyperref[exo-16-10]{16.10}, \hyperref[exo-16-11]{16.11}, \hyperref[exo-16-12]{16.12}, \hyperref[exo-16-14]{16.14}\\[2mm]
\bottomrule\end{tabularx}\par\vspace{4mm}
\otIntertitre{17. Espace, pyramides et cônes}\begin{tabularx}{\linewidth}{@{}Yl@{}}\toprule
Pyramides, cônes et perspective & \hyperref[exo-17-1]{17.1}, \hyperref[exo-17-2]{17.2}, \hyperref[exo-17-3]{17.3}, \hyperref[exo-17-4]{17.4}\\[2mm]
Repérage dans un pavé droit & \hyperref[exo-17-5]{17.5}, \hyperref[exo-17-6]{17.6}\\[2mm]
Patrons de pyramides et de cônes & \hyperref[exo-17-9]{17.9}, \hyperref[exo-17-10]{17.10}\\[2mm]
Volumes, unités et assemblages & \hyperref[exo-17-7]{17.7}, \hyperref[exo-17-8]{17.8}, \hyperref[exo-17-11]{17.11}, \hyperref[exo-17-12]{17.12}, \hyperref[exo-17-13]{17.13}, \hyperref[exo-17-14]{17.14}, \hyperref[exo-17-15]{17.15}\\[2mm]
\bottomrule\end{tabularx}\par\vspace{4mm}
\otIntertitre{18. Algorithmique et programmation par blocs}\begin{tabularx}{\linewidth}{@{}Yl@{}}\toprule
Séquences, entrées et variables & \hyperref[exo-18-1]{18.1}, \hyperref[exo-18-2]{18.2}, \hyperref[exo-18-3]{18.3}, \hyperref[exo-18-5]{18.5}, \hyperref[exo-18-6]{18.6}\\[2mm]
Conditions et cas limites & \hyperref[exo-18-4]{18.4}, \hyperref[exo-18-9]{18.9}, \hyperref[exo-18-15]{18.15}\\[2mm]
Boucles et constructions géométriques & \hyperref[exo-18-7]{18.7}, \hyperref[exo-18-8]{18.8}, \hyperref[exo-18-11]{18.11}, \hyperref[exo-18-13]{18.13}\\[2mm]
Événements, interaction et mise au point & \hyperref[exo-18-10]{18.10}, \hyperref[exo-18-12]{18.12}, \hyperref[exo-18-14]{18.14}\\[2mm]
\bottomrule\end{tabularx}\par\vspace{4mm}
\clearpage
